\section{Running the experiments}
\label{app:scripts}

All experiments of these notes are in the folder \texttt{code/}. They need Python with the packages
\texttt{torch}, \texttt{transformer\_lens}, \texttt{datasets} and \texttt{matplotlib}, and no GPU. Each script
writes its numbers to a file \texttt{code/results\_*.json} and its figures to \texttt{figures/}.
\Cref{tab:scripts} lists them with their running time on a laptop with four CPU cores. The command
\texttt{sh code/run\_all.sh} runs all of them one after the other.

\begin{table}[h]
\centering\small
\begin{tabular}{@{}llp{0.48\textwidth}r@{}}
\toprule
Script & Sessions & What it does & Minutes \\
\midrule
\texttt{exp\_basics.py} & 1, 2 & tokens, predictions, attention patterns & 1 \\
\texttt{exp\_ioi.py} & 1, 2, 6 & the running example, logit lens, direct logit attribution, patching of the stream and of all heads & 2 \\
\texttt{exp\_directions.py} & 3 & probes for the gender of a name, steering & 2 \\
\texttt{exp\_neurons.py} & 4 & strongest activations of single neurons & 1 \\
\texttt{exp\_superposition.py} & 4 & trains the toy model, random directions & 5 \\
\texttt{exp\_sae.py} & 5 & downloads a public sparse autoencoder and evaluates it & 3 \\
\texttt{exp\_induction.py} & 7 & repeated random tokens, induction scores, ablation & 1 \\
\texttt{exp\_modadd.py} & 8 & trains a network on addition modulo 113 and analyses it & 15 \\
\bottomrule
\end{tabular}
\caption{The scripts that produce all numbers and data figures of these notes.}
\label{tab:scripts}
\end{table}

The scripts use the library TransformerLens \citep{nanda2022transformerlens}. This library gives access to every
intermediate activation of a model through named \emph{hooks}. The listing below is a complete patching experiment
in thirty lines. It repeats \cref{exp:residpatch} for a single pair of prompts, and it is a good starting point for
a presenter who wants to try a variation. Two related tools are nnsight \citep{fiottokaufman2025nnsight}, which
works with models that are too large for a laptop, and the website Neuronpedia \citep{lin2023neuronpedia}, which
shows the latents of published sparse autoencoders.

\lstinputlisting[style=tutpython]{code/minimal_patching.py}

The listing was tested with version 4.4 of the library. Older tutorials use a class called
\texttt{HookedTransformer}, which was removed in version 4.0. These tutorials still run with an older version of
the library, which can be installed with \texttt{pip install "transformer-lens<4"}.

\section{Glossary}
\label{app:glossary}

{\small \setlength{\columnsep}{20pt}
\begin{multicols}{2}
\begin{description}[style=sameline, leftmargin=0pt, itemsep=1.5pt, font=\normalfont\bfseries]
\item[Ablation.] Replacing the value of a node by zero, by its mean or by its value on another prompt, and
measuring the change of the output.
\item[Activation.] Any intermediate value that the model computes on a given input.
\item[Activation patching.] Running the model on one prompt while overwriting one node with its value from a
second, minimally different prompt.
\item[Attention head.] The component that moves information between positions.
\item[Attention pattern.] The weights with which a head at a destination reads from the sources.
\item[Autoencoder.] A network that is trained to reproduce its own input after passing it through an
intermediate representation.
\item[Circuit.] A subgraph of the model's components together with a description of what each one computes on
a task.
\item[Clean and corrupted prompt.] The two prompts of a patching experiment. They differ only in the aspect
under study.
\item[Dictionary direction.] A column of the decoder of a sparse autoencoder.
\item[Difference of means.] The average activation on inputs with a concept minus the average on inputs
without it.
\item[Direct logit attribution.] The term that one component contributes to a logit through its own addition
to the residual stream.
\item[Embedding.] The vector with which a token enters the model.
\item[Faithful.] Said of an explanation that correctly describes what the model computes. Opposed to merely
plausible.
\item[Feature.] A property of the input or of the intermediate computation that the model represents.
\item[Forward and backward pass.] Computing the output of the model for an input, and computing the gradient of
a number at the output with respect to all weights and activations.
\item[Grokking.] Generalisation that sets in long after the training data have been fitted.
\item[Induction head.] A head that attends to the token after an earlier occurrence of the current token.
\item[Latent.] One unit of a sparse autoencoder. A candidate for a feature.
\item[Linear probe.] A logistic regression that predicts a label from an activation.
\item[Logit.] The score of a token before the softmax.
\item[Logit difference.] The logit of the correct token minus that of a chosen alternative.
\item[Logit lens.] Applying the final read-out to an intermediate residual stream.
\item[Neuron.] One coordinate of the hidden layer of an MLP.
\item[Polysemantic.] Said of a neuron that responds to several unrelated features.
\item[Prompt.] The text that is given to the model as input.
\item[Residual stream.] The vector that is kept for each position and to which all components add their
output.
\item[Sparse autoencoder.] A wide one-layer autoencoder trained to reconstruct activations with few active
latents.
\item[Steering.] Adding a direction to the residual stream to change a behaviour.
\item[Superposition.] The storage of more features than dimensions along almost orthogonal directions.
\item[Token.] A word or word fragment. The unit in which the model reads and writes text.
\item[Unembedding.] The vectors that turn the final residual stream into logits.
\item[Weight decay.] A part of the training step that shrinks every weight by a small fraction of its value.
\end{description}
\end{multicols}
}
